\chapter{Arranque-pare sencillo y reversible.}


\section{\obj}
\capacidad
\begin{itemize}
    \item Estudiar  las representaciones simbólicas normalizadas de los elementos de control eléctrico, según norma NEMA ICS 19-2002 y IEC 60617.
    \item Seleccionar contactores y relevadores térmicos.
    \item Estudiar la representación gráfica del estandar NEMA ICS 19-2002 o IEC 60617. 
    \item Diseñar e implementar circuitos de control para arranque sencillo de un motor trifásico.
    \item Diseñar e implementar circuitos de control para arranque-pare reversible de un  motor trifásico.
\end{itemize} 

 
\section{Equipos y materiales}
Para este laboratorio de necesitaran:
\begin{itemize}
    \item 1 motor trifásico 1/2 hp de 3 puntas
	\item 1 relevador térmico
	\item 1 botón pulsar cerrado
	\item 2 botones pulsadores abiertos
	\item 1 multimetro.
\end{itemize}

\section{Metodología}

Este laboratorio tiene una duración de 4 lecciones, repartidas en dos semanas. Los estudiantes deben mostrar durante las clases programadas las tres actividades propuestas. Deben recabar fotografías y resultados de los equipos de medición para elaborar las evidencias. Las evidencias se subirán al TecDigital la semana siguiente finalizadas las actividades.

\section{Práctica en Clase}

\subsection{Actividad 1}
  Los ingenieros Electricistas, electromecánicos y de Mantenimiento pueden diseñar sus circuitos de control eléctricos, sin embargo cada fabricante posee libros de diagramas estandarizados como el que brinda  Schneider con su documento \href{https://download.schneider-electric.com/files?p_enDocType=Catalog&p_File_Name=0140CT9201.pdf&p_Doc_Ref=0140CT9201&_ga=2.160182845.1491407618.1677858387-1733391740.1677858387}{\textit{Wiring Diagram Book}}\cite{Scheneider2}.
 
 Estudie toda la simbología NEMA, Tabla 1 y Tabla 2 de la referencia \cite{Scheneider2} y asocie los con el elementos físicos y su funcionamiento.
 Deduzca la ecuación lógica del circuito eléctrico de control, para esto use un mapa de Karnaugh o la ecuacion caracteristica del cerrojo S-R prioridad al reset.
 Realice el alambrado del circuito de control y luego del circuito de potencia del arranque sencillo de un motor trifásico de 1/2 hp.

\begin{equation*}
    M^+ = \overline{P} \cdot \overline{Sc} \cdot (A+M)
\end{equation*}

\begin{figure}[H]
    \centering
    \begin{tikzpicture}[circuit plc ladder,scale=1.5]
    \draw(0,0)  
    to [ncpb,l=$P$] ++ (1,0) coordinate(M1)
    to [nopb,l=$A$] ++ (1,0) coordinate(M2)
    to [esource, l=$M$] ++ (1,0)
    to [contact NC={info={$Sc$}}] ++ (1,0) coordinate(laddertopright)
    (M1) -- ++ (0,-0.5) 
    to [contact NO={info={$M$}}] ++ (1,0) -- (M2)
    ;
    \ladderrungend{1}
    \ladderpowerrails
    \end{tikzpicture}
    \caption{Circuito de control: arranque-pare secillo}
\end{figure}

\begin{figure}[H]
\centering
    \begin{circuitikz}[cute inductors,scale=0.7]
    \node[esourceshape,name=motor] at (8,-2){};
    \draw
    (0,0) node[left]{$L_1$} 
    --++ (1.5,0) coordinate(L11) -- ++ (0.5,0)
    to [C,l=$M$] 
    ++ (2,0) to [wfuse,l=$S_{c}$] ++ (2,0) -- (motor.120)
    (0,-2) node[left]{$L_2$} 
    --++ (1,0) coordinate(L21) -- ++ (1,0)
    to [C,l=$M$] 
    ++ (2,0) to [wfuse,l=$S_{c}$] ++ (2,0) -- (motor.180)
    (0,-4) node[left]{$L_3$} 
    --++ (0.5,0) coordinate(L31) -- ++ (1.5,0)
    to [C,l=$M$] 
    ++ (2,0) to [wfuse,l=$S_{c}$] ++ (2,0) -- (motor.240)
    (motor.center) node{$M$}
    ; 
    
    \end{circuitikz}
    \caption{Circuito de potencia: arranque-pare secillo}
    \label{fig:potencia-secuencial}
\end{figure}


\subsubsection{Conteste las preguntas:}

¿Es el circuito alambrado control a dos o tres hilos? Explique. ¿El circuito de control funcionó según la tabla de verdad?, Que pasa si se pierde una fase, como se comporta el circuito?. Que diferencia existe entre un arrancador automático y uno manual en el caso que exista un corte de flujo eléctrico? Explique.
 Hasta ahora los circuitos fueron representados según la norma NEMA, muestre el circuito de control y potencia según la simbología de la norma IEC.

\subsection{Actividad 2}

Deduzca el circuito de control de un arranque-pare-reversible. Realice el alambrado del circuito de control y luego del circuito de potencia del circuito para un trifásico de 1/2 hp.


\begin{align*}
    F^+ &= \overline{P} \cdot \overline{R} \cdot \overline{Sc} \cdot (A_F + F) \\
    R^+ &= \overline{P} \cdot \overline{F} \cdot \overline{Sc} \cdot (A_R + R)
\end{align*}

\begin{figure}[H]
    \centering
    \begin{tikzpicture}[circuit plc ladder,scale=1.5]
    \draw(0,0)  
    to [ncpb,l=$P$] ++ (1,0) coordinate(AF1)
    to [nopb,l=$A_F$] ++ (1,0) coordinate(AF2)
    to [contact NC={info={$R$}}] ++ (1,0)
    to [esource, l=$F$] ++ (1,0)
    to [contact NC={info={$Sc$}}] ++ (1,0) coordinate(laddertopright)
    (AF1) -- ++ (0,-0.5) 
    to [contact NO={info={$F$}}] ++ (1,0) -- (AF2)
    (M1) -- ++ (0,-2) coordinate(AR1)
    to [nopb,l=$A_R$] ++ (1,0) coordinate(AR2)
    to [contact NC={info={$F$}}] ++ (1,0)
    to [esource, l=$R$] ++ (1,0) -- ++ (0,2)
    (AR1) -- ++ (0,-0.5) 
    to [contact NO={info={$R$}}] ++ (1,0) -- (AR2)
    ;
    \ladderrungend{1}
    \ladderpowerrails
    \end{tikzpicture}
    \caption{Circuito de control: arranque Estrella-Delta reversible}
\end{figure}

\begin{figure}[H]
    \centering
    \begin{circuitikz}[cute inductors,scale=0.7]
    \node[esourceshape,name=motor] at (10,-2){};
    \draw
    (0,0) node[left]{$L_1$} 
    --++ (1.5,0) coordinate(L11) -- ++ (0.5,0)
    to [C,l=$F$] 
    ++ (2,0) -- ++ (1.5,0) coordinate(L12) -- ++ (0.5,0) coordinate(L13)
    (0,-2) node[left]{$L_2$} 
    --++ (1,0) coordinate(L21) -- ++ (1,0)
    to [C,l=$F$] 
    ++ (2,0) -- ++ (1,0) coordinate(L22) -- ++ (1,0) coordinate(L23) 
    (0,-4) node[left]{$L_3$} 
    --++ (0.5,0) coordinate(L31) -- ++ (1.5,0)
    to [C,l=$F$] 
    ++ (2,0) -- ++ (0.5,0) coordinate(L32) -- ++ (1.5,0) coordinate(L33)
    (L31) 
    to[short,*-] 
    ++ (0,6) -- ++ (1.5,0) 
    to [C,l=$R$] 
    ++ (2,0) -- ++ (0.5,0) 
    to[short,-*]
    (L32)
    (L21) 
    to[short,*-] 
    ++ (0,6) -- ++ (1,0) 
    to [C,l=$R$] 
    ++ (2,0) -- ++ (1.5,0) 
    to[short,-*]
    (L12)
    (L11) 
    to[short,*-] 
    ++ (0,6) -- ++ (0.5,0) 
    to [C,l=$R$] 
    ++ (2,0) -- ++ (1,0) 
    to[short,-*]
    (L22)
    (L13) to[wfuse,l=$S_{c}$] ++ (2,0) -- (motor.120)
    (L23) to[wfuse,l=$S_{c}$] ++ (2,0) -- (motor.180)
    (L33) to[wfuse,l=$S_{c}$] ++ (2,0) -- (motor.240)
    (motor.center) node{$M$}
    ;
    \end{circuitikz}
    \caption{Circuito de potencia: arranque Estrella-Delta reversible}
    \label{fig:potencia-estrella-delta}
\end{figure}

\subsubsection{Conteste las preguntas:} 

 ¿El circuito de control funcionó según la tabla de verdad?, ¿Que pasa si se presionan los botones pulsadores de arranque hacia la derecha y arranque hacia la izquierda al mismo tiempo? Explique usando el concepto de enclavamiento. 
 Muestre el circuito de control y potencia según la simbología de la norma IEC.



